% \section{Background}
\section{Preliminary}\label{sec-pre}
This section introduces the key concepts and definitions related to the TKGQA task and its formulation.

\noindent\textbf{Temporal Knowledge Graph.}  Formally defined as $\mathcal{G} =  (\mathcal{E}, \mathcal{R}, \mathcal{T}, \mathcal{F})$, where $\mathcal{E}$, $\mathcal{R}$, $\mathcal{T}$, and $\mathcal{F}$ represent the sets of entities, relations, timestamps, and facts, respectively~\cite{caiSurveyTemporalKnowledge2024}. Each temporal fact \( f \in \mathcal{F} \) encodes relationships between entities with explicit timestamps, typically represented as $(s,r,o, \tau)$, quadruples where $s,o\in \mathcal{E}, r\in \mathcal{R},$ and $\tau\in\mathcal{T}$.

TKGs support multiple representation forms of temporal facts to capture temporal information effectively:
\begin{itemize}
\item \textbf{Compound Value Types (CVTs)}: In Freebase~\cite{bollackerFreebaseCollaborativelyCreated2008}, CVTs are attributes that consist of multiple related values. For example, the \texttt{marriage} predicate is a CVT with attributes like \texttt{marriage.spouse} and \texttt{marriage.date}.
\item \textbf{Timestamped Triples}: Triples with temporal annotations, e.g., \texttt{<Malia Obama, date\_of\_birth, 1998-07-04>}
\item \textbf{N-array Tuples}: Encode multi-attribute facts through qualifiers, e.g., \texttt{<Barack Obama, position held, President of the United States; start date, 20-01-2009; end date, 20-01-2017>}
\item \textbf{Quintuples}: Compact representation of interval-based facts, e.g., \texttt{<Barack Obama, position held, President of the United States, 2009, 2017>}. 
\item \textbf{Events}: Encode significant events such as \texttt{<WWII, significant event, occurred, 1939, 1945>}~\cite{saxenaQuestionAnsweringTemporal2021}.
\end{itemize}

Notably, some representations can be decomposed into standard quadruples. For instance, the presidential term quintuple \texttt{<Barack Obama, position held, President of the United States, 2009, 2017>} can be decomposed into annual quadruples: \texttt{<Barack Obama, position held, President of the United States, 2009>}, \texttt{<…, 2010>}, …, \texttt{<…, 2017>}. Such normalization of time-based facts ensures that the system can be applied across various TKGs.

\noindent\textbf{Temporal Question.}
 A temporal question contains at least one temporal constraint or requires timestamps as its answer~\cite{jiaTempQuestionsBenchmarkTemporal2018}. 
 
The \textit{temporal constraint}, serving as the key distinguishing text span in temporal questions, specifies a condition tied to a particular time point or interval that the answer and its supporting evidence must satisfy. This constraint typically combines a \textit{temporal expression} (e.g., 2023) and a \textit{temporal signal} (e.g., after) that co-occur as contiguous spans within the question text.
 
To dissect this structure further, \textit{temporal expressions} refer to time points or intervals in temporal questions, which can vary in granularity (e.g., May 11th, 2024)~\cite{pustejovskyTimeMLRobustSpecification,huangDomainSensitiveTemporalTagging2018}. Complementing temporal expressions, \textit{temporal signals} indicate the relationships between them and act as trigger words, imposing constraints on the answer or supporting evidence (e.g., in, after, or during).

\noindent\textbf{Temporal Knowledge Graph Question Answering.} Given a TKG $\mathcal{G}$ and a temporal question $q$ expressed in natural language, the task of TKGQA aims to infer the correct answer from the structured data. The answers consist of entity sets $\left\{e \mid e \in \mathcal{E}\right\}$ or timestamps $\left\{\tau \mid \tau \in \mathcal{T}\right\}$ within $\mathcal{G}$.

\noindent\textbf{TKGQA Datasets.} A TKGQA dataset $\mathcal{D}$ comprises pairs of questions and answers $(q, a)$. Different datasets may include additional information, such as entity identifiers~\cite{neelamSYGMASystemGeneralizable2021}, temporal expressions within $q$, or question types. 

\begin{table}[]
\caption{Knowledge representation form of TKGQA datasets.\label{tab:knowledge}}
\begin{tabular}{lll}
\toprule
\textbf{Dataset}       & \textbf{Underlying TKG}        & \textbf{\makecell{Knowledge Representation \\ Form}}                                                   \\ \midrule
TempQuestions & Freebase   & CVTs                                                                            \\ \midrule
TimeQuestions & Wikidata   & \begin{tabular}[c]{@{}l@{}}Timestamped Triples and\\ N-array Tuple\end{tabular} \\ \midrule
CronQuestions & Wikidata   & Quintuples or Quadruples                                                        \\ \midrule
MultiTQ       & ICEWS05-15 & Quadruples \\ \bottomrule                
\end{tabular}
\end{table}

Since different datasets are constructed from different knowledge graphs, their representation form of temporal facts also varies. We list a few examples in Table~\ref {tab:knowledge}. 
% For the underlying TKG, some datasets directly provide a processed TKG, while others indicate the source TKG related to the question through annotations~\cite{jiaComplexTemporalQuestion2021}.
% Detailed descriptions of these datasets can be found in Appendix~\ref{sec-datasets}.
